\chapter{Списки (Lists)}\label{lst}

Списки "--- це структура даних, що, мабуть, найчастіше використовується в мові \Python. Вони надають програмісту широкий спектр можливостей для гнучкого і ефективного вирішення найрізноманітніших прикладних завдань. Застосування списків дозволяє отримати елегантне і коротке Розв'язання, економлячи десятки рядків коду і час розробки. За своєю потужністю, гнучкістю та простотою використання списки перевершують аналогічні структури (масиви, динамічні масиви) в інших мовах програмування, що ж час переконатися в цьому на власному досвіді.

\section{Створення списків}

\begin{defn}{Список "--- це}
\textbf{змінна послідовність об'єктів довільних типів.}
\end{defn}
\noindent Як бачимо, принциповою відмінністю списків від раніше розглянутих нами структур даних є їх \textbf{змінюваність}. Саме ця властивість робить списки настільки корисними та ефективними. Зокрема, це означає, що
\begin{itemize}
\item список має змінну довжину;
\item елементам списку можна надавати нові значення;
\item можна додавати та видаляти елементи списку.
\end{itemize}

Список у коді програми можна створити за допомогою перерахування його елементів у квадратних дужках:
\begin{py}{Приклади: створення списків за допомогою літералу}
L = [] # порожній список
L = [1, 2, 3] # список із трьох елементів
L = [' abc ', (0, 1), [' def ', 2, 3]] # список із трьох елементів
\end{py}
\noindent З прикладів бачимо, що елементами списку можуть бути числа, рядки, кортежі та інші списки. Цього виявляється більш ніж достатньо для вирішення багатьох завдань.

Як і кортеж, \textbf{список можна створити з будь-якого об'єкта, що ітерується}:
\begin{py}{Приклади: створення списку з об'єкта, що ітерується}
L = list (range (10)) # список з віртуальної послідовності
L = list (' text ') # список з рядка
L = list ((1, 2, 3)) # список з кортежу
\end{py}

\subsection{Конструктори списків} \label{lstcomp}

Зручним інструментом для створення списку є \textbf{конструктор (генератор) списку} (list comprehension). Він дозволяє створити список на підставі \textbf{деякого загального правила}, як це прийнято при завданні множин у математиці. Тут ми використовуємо іншу послідовність, задаючи щось подібне до формули для обчислення елементів списку. Наприклад, можна створити список квадратів перших 100 натуральних чисел таким чином:
\begin{py}{Приклад: конструктор списку}
L = [x**2 for x in range (1, 101)]
\end{py}
\noindent Можна задати \textbf{додаткова умова відбору}:
\begin{py}{Приклад: конструктор списку з умовою відбору}
L = [x**2 for x in range (1, 101) if x % 2 == 1]
\end{py}
\noindent Отримаємо список квадратів лише непарних чисел від 1 до 100.

\medskip

Наведемо ще кілька прикладів створення списку за допомогою конструктора.
\begin{py}{Приклади: конструктори списків}
L = ['a'*k for k in range (5)] # ['', 'a', 'aa', 'aaa', 'aaaa']
L = [i+1 for i in range (5)] # [1, 2, 3, 4, 5]
L = [x for x in range (1, 101) if x%5==0 and x%7==0] # [35, 70]
\end{py}

Способи введення списків з клавіатури та файлу будуть розглянуті \hyperref [listinput]{нижче}.

\section{Основні функції та методи}

Для початку зауважимо, що до списків застосовуються всі операції, розглянуті нами для кортежів:
\begin{itemize}
\item довжина списку: \mi{len (L)}
\item конкатенація: \mi{L 1+ L 2}
\item повторення: \mi{L * n} або \mi{n * L}
\item доступ за індексом: \mi{L [i]}
\item доступ за індексом до вкладених елементів: \mi{L [i] [j]}, \mi{L [i] [j] [k]} і т.д.
\item зрізи: \mi{L[i:j:step]}
\item перевірка на входження: \mi{elem in L} (\mi{True}/\mi{False})
\item підрахунок кількості входжень елемента: \mi{L. count (elem)}
\item індекс першого входження елемента: \mi{L. index (elem)}
\end{itemize}

\medskip

Далі розглянемо специфічні операції, які застосовуються лише до списків.

\subsection{Надання нових значень}

\begin{pydef}
L [i] = new _ elem # зміна значення елемента [i]
L [i:j] = new _ list # привласнення зрізу
L [i:j:step] = new _ list # привласнення зрізу з кроком
\end{pydef}

Надавати нові значення можна як окремим елементам, так і зрізам.

Елементу списку \mi{L[i]} можна присвоїти будь-який об'єкт "--- число, рядок, кортеж, список. При цьому \textbf{зміни вносяться до самого вихідного списку}; новий список не створюється.

У разі присвоєння зрізу є обмеження: об'єкт праворуч від оператора присвоєння має бути ітерованим (рядок, кортеж, список, віртуальна послідовність). Об'єкт зліва "--- зріз елементів списку "--- завжди є списком.

У разі присвоєння зрізу з кроком обмеження ще більш жорстке: довжини об'єктів зліва і праворуч від оператора присвоєння повинні збігатися.

Розглянемо з прикладу, як це працює.

\begin{pyr}{}
\begin{py}{Приклади: надання нових значень елементам списку}
L1 = list(range(10))
print('First list:', L1)
L1[3] = 'bug'
print('Modified L1. Step 1:', L1)
L1[5:] = [x for x in range (3)]
print('Modified L1. Step 2:', L1)
L1[-1] = [x for x in range (3)]
print('Modified L1. Step 3:', L1)
L2 = list('abrakadabra')
print('Second list:', L2)
L2[1:len(L2)-1:2] = '!'*5
print('Modified L2:', L2)
\end{py}
\tcblower
\begin{py}{}
>>>
First list: [0, 1, 2, 3, 4, 5, 6, 7, 8, 9]
Modified L1. Step 1: [0, 1, 2, 'bug', 4, 5, 6, 7, 8, 9]
Modified L1. Step 2: [0, 1, 2, 'bug', 4, 0, 1, 2]
Modified L1. Step 3: [0, 1, 2, 'bug', 4, 0, 1, [0, 1, 2]]
Second list: ['a', 'b', 'r', 'a', 'k', 'a', 'd', 'a', 'b', 'r', 'a']
Modified L2: ['a', '!', 'r', '!', 'k', '!', 'd', '!', 'b', '!', 'a']
>>>
\end{py}
\end{pyr}

Зверніть увагу на відмінності в результатах надання списку-генератора \- \mi{[x for x in range (3)]} зрізу, а потім окремому елементу. У першому випадку весь зріз \mi{L1[5:]} (це список довжини 5) був замінений новим списком (довжини 3), при цьому елементи нового списку додаються у вихідний як самостійні об'єкти. У другому випадку один елемент списку замінено списком довжини 3.

\subsection{Видалення елементів за індексом}

\begin{pydef}
del L [i] # видалення елемента [i]
del L [i : j : step] # видалення зрізу
\end{pydef}

Видаляти зі списку можна також окремі його елементи, так і зрізи.

\begin{pyr}{}
\begin{py}{Приклади: видалення елементів списку}
L = list(range(10))
print('L:', L)
del L[0]
print('Modified L. Step 1:', L)
del L[1:4]
print('Modified L. Step 2:', L)
del L[::2]
print('Modified L. Step 3:', L)
\end{py}
\tcblower
\begin{py}{}
>>>
L: [0, 1, 2, 3, 4, 5, 6, 7, 8, 9]
Modified L. Step 1: [1, 2, 3, 4, 5, 6, 7, 8, 9]
Modified L. Step 2: [1, 5, 6, 7, 8, 9]
Modified L. Step 3: [5, 7, 9]
>>>
\end{py}
\end{pyr}
\noindent Результати не вимагають особливих коментарів. Зауважимо лише, що команда \mi{del} <<бачить>> лише адреси (індекси) об'єктів, що видаляються. Зрозуміло, було б добре мати можливість видаляти елементи, вказуючи їх значення, а чи не адреси. Про це "--- у наступному параграфі.

\subsection{Видалення елементів за значенням} \label{remove}

\begin{pydef}
L.remove(elem)
\end{pydef}

Даний метод здійснює видалення (removing) першого (ліворуч) входження елемента \mi{elem} до списку \mi{L}. Якщо елемент у списку не знайдено, видається повідомлення про помилку.

Реалізувати видалення \textbf{всіх входжень} елемента в список можна за допомогою наступного \mi{while}-циклу:
\begin{py}{}
while elem in L :
L.remove(elem)
\end{py}

Звернімо увагу на дещо незвичний вид запису команди \mi{remove()}. Вона записана окремим рядком програми, при цьому результат начебто нікуди не запам'ятовується. Справа в тому, що результатом виконання команди є \textbf{зміна самого списку \mi{L}}; ніякий об'єкт  не повертається\footnote{Якщо бути точним, даний метод повертає об'єкт типу \sf{NoneType}, тобто <<нічого>>.}.

\begin{pyr}{}
\begin{py}{Приклад: видалення елементів списку за значенням}
L = [1, 2, '1', [1, 2, 3]]
L.remove(1)
print('L:', L)
\end{py}
\tcblower
\begin{py}{}
>>>
L: [2, '1', [1, 2, 3]]
>>>
\end{py}
\end{pyr}
\noindent Подумайте над тим, чому не було видалено елемент \mi{'1'}, а також одиниця, що входить до списку \mi{[1, 2, 3]}.

%Зауважимо, що елемент \mi{'1'} не був видалений, тому що це рядок, а ми видаляли число 1. Також не була видалена одиниця, що входить до списку \mi{[1, 2, 3]}, як об'єкт , що не є елементом списку \mi{L}.

\subsection{Вилучення елементів за індексом} \label{pop}

\begin{pydef}
L.pop(i)
L.pop()
\end{pydef}

Метод видаляє зі списку \mi{L} елемент з індексом \mi{[i]}, при цьому повертається значення цього елемента \mi{L[i]}. Таким чином, ми можемо дізнатися, який саме об'єкт було видалено. Якщо не вказати індекс \mi{[i]}, буде вилучено останній елемент списку.

\begin{pyr}{}
\begin{py}{Приклади: вилучення елементів списку}
L = [1, 2, '1', [1, 2, 3]]
what = L.pop(3)
print('Removing:', what)
print('New list:', L)
print('Removing:', L.pop())
print('New list:', L)
\end{py}
\tcblower
\begin{py}{}
>>>
Removing: [1, 2, 3]
New list: [1, 2, '1']
Removing: 1
New list: [1, 2]
>>>
\end{py}
\end{pyr}

\subsection{Додавання елементів}

\begin{pydef}
L.append(obj)
\end{pydef}

Метод \mi{append()} додає об'єкт \mi{obj} в кінець списку \mi{L} (appending "--- додавання). При цьому об'єкт \mi{obj} може бути будь-якого типу.

\begin{pydef}
L.extend(iter_obj)
\end{pydef}

Метод \mi{extend()} додає \textbf{об'єкт, що ітерується} \mi{iter_obj} в кінець списку \mi{L}, при цьому додавання відбувається \textbf{поелементно} "--- т. н. <<розширення списку >> (extension; звідси назва методу). Різниця між \mi{append()} і \mi{extend()} буде зрозуміла на прикладах.

\begin{pydef}
L.insert(i, obj)
\end{pydef}

Метод \mi{insert()} здійснює вставку (insertion) об'єкта \mi{obj} в позицію \mi{L[i]} списку \mi{L}. Якщо значення \mi{i} перевищує \mi{len(L)}, то вставка здійснюється в кінець списку.

Як і \hyperref[remove]{метод \mi{remove()}}, методи \mi{append()}, \mi{extend()} і \mi{insert()} роблять зміну списку \mi{L} і жодного об'єкта не повертають.

\begin{pyr}{}
\begin{py}{Приклади: додавання елементів до списку}
L = [0, 1]
L.append('abc')
print('Step 1:', L)
L.extend('abc')
print('Step 2:', L)
L.insert(0, -1)
print('Step 3:', L)
\end{py}
\tcblower
\begin{py}{}
>>>
Step 1: [0, 1, 'abc']
Step 2: [0, 1, 'abc', 'a', 'b', 'c']
Step 3: [-1, 0, 1, 'abc', 'a', 'b', 'c']
>>>
\end{py}
\end{pyr}
\noindent Зверніть увагу на те, як спрацювали методи \mi{append()} та \mi{extend()} з одним і тим же аргументом "--- рядком \mi{'abc'}. Зауважимо, що при додаванні неітерованого об'єкт метод \mi{extend()} повертає помилку.

\subsection{Інше}

\begin{pydef}
L.reverse()
\end{pydef}

Метод \mi{reverse()} змінює порядок проходження елементів списку на зворотний. Результатом є зміна вихідного списку \mi{L}.

\begin{pydef}
L.sort ()
\end{pydef}

Метод \mi{sort()} проводить сортування списку за зростанням елементів. Якщо список містить різнорідні елементи, сортування відбувається в так званому «словниковому» порядку. За подробицями радимо звернутись до документації. Зауважимо, що такий метод сортування не завжди достатній для вирішення завдань, тому далі нами будуть розглянуті декілька стандартних \hyperref [sort]{алгоритмів сортування}.

Подумайте, як за допомогою комбінації методів \mi{reverse()} та \mi{sort()} можна відсортувати список у порядку спадання.

\begin{pydef}
L1 = sorted(L)
L1 = sorted(L, reverse = True)
\end{pydef}

Функція \mi{sorted()} повертає відсортований список. За замовчуванням сортується за зростанням. Параметр \mi{reverse = True} дає сортування за спаданням. Відмінність від методу \mi{sort()} у тому, що в даному випадку створюється новий об'єкт-список (початковий список \mi{L} не змінюється), тоді як \mi{sort()} виконує сортування <<на місці> >, змінюючи сам список \mi{L}.

У модулі \mi{random} визначено метод
\begin{pydef}
random.choice(L)
\end{pydef}
\noindent для вибірки випадкового елемента з непустого списку \mi{L}\footnote{Замість списку може використовуватись будь-яка послідовність.}. Якщо список \mi{L} порожній, повертається помилка.


\section{Способи введення списків}\label{listinput}

Звернімося до створення списків у коді програми. Вище нами було розглянуто два способи:
\begin{itemize}
\item перерахування елементів;
\item конструктор списку.
\end{itemize}
У даному параграфі розглянемо введення списків \hyperref [inkeyb]{з клавіатури}, \hyperref [infile]{з файлу}, а також \hyperref [inrand]{заповнення випадковими числами}.

\subsection{Введення з клавіатури}\label{inkeyb}

Нехай потрібно ввести цілий список з \mi{n} елементів.

\cleardoublepage

\paragraph{Поелементне введення}\mbox{}

\begin{defn}{Поелементне введення списку з клавіатури}
\begin{py}{}
n = int (input ('Введіть кількість елементів списку:'))
L = [] # створюємо порожній список
print ('Введіть елементи:')
for i in range (n) : цикл з n кроків
    x = int (input ()) # вводимо черговий елемент і натискаємо ENTER
    L. _ append (x) # додаємо елемент у кінець списку
\end{py}
\end{defn}
\noindent Ці ж дії можна записати коротше, якщо використовувати конструктор списку:
\begin{defn}{Поелементне введення списку: використання конструктора}
\begin{py}{}
n = int (input ('Введіть кількість елементів списку:'))
print('Введіть елементи:')
L = [int(input()) for i in range (n)]
\end{py}
\end{defn}

\paragraph{Введення елементів в один рядок}
Якщо елементів не дуже багато, можна ввести їх через пробіл в один рядок, а потім перетворити рядок на список:
\begin{defn}{Введення елементів списку з клавіатури в один рядок}
\begin{py}{}
L = []
print ('Введіть елементи списку через пробіл:')
s = input () # вводимо рядок: елементи поділяються пробілами
L = [int(x) for x in s.split()] # формуємо список
\end{py}
\end{defn}

\subsection{Введення з файлу}\label{infile}

Нехай у файлі \mi{input.txt} зберігаються цілі числа по одному в рядку. Сформуємо список із цих чисел.

\begin{defn}{Введення елементів списку з файлу}
\begin{py}{}
f = open('input.txt', 'r')
L = []
for line in f :
    L. _ append (int (line [:-1])) # в даному випадку можна і int (line)
f.close()
\end{py}
\end{defn}
\noindent Якщо використовувати конструктор, код стане коротшим:
\begin{defn}{Введення елементів списку з файлу: використання конструктора}
\begin{py}{}
f = open('input.txt', 'r')
L = [int(line[:-1]) for line in f]
f.close()
\end{py}
\end{defn}

\subsection{Заповнення випадковими числами} \label{inrand}

Створимо список з \mi{n} випадкових цілих чисел з інтервалу \mi{(-10,10)}.

\begin{defn}{Заповнення списку випадковими числами}
\begin{py}{}
from random import *
n = int (input ('Кількість елементів списку:'))
L = []
for i in range (n):
    L.append(randint(-10, 10))
\end{py}
\end{defn}
\noindent З конструктором:
\begin{defn}{Заповнення списку випадковими числами: використання конструктора}
\begin{py}{}
from random import *
n = int (input ('Кількість елементів списку:'))
L = [randint(-10, 10) for i in range (n)]
\end{py}
\end{defn}


\section{Виведення списків на екран}

Стандартна функція \mi{print(L)} виводить список у вигляді послідовності елементів через кому, укладену в квадратні дужки. А якщо нам потрібно вивести список якось інакше?

Наприклад, виведемо список без квадратних дужок, розділяючи елементи точкою із комою і пробілом. Для цього стане в нагоді наш старий знайомий "--- метод \mi{join()}.

\begin{defn}{<<Красивий>> виведення списку}
\begin{py}{}
print('; '.join([str(x) for x in L])))
\end{py}
\end{defn}

Зрозуміло, можна використовувати інші роздільники.


\section{Алгоритми сортування списків}\label{sort}

У цьому параграфі розглянемо три найпростіші способи сортування елементів одновимірного списку: \hyperref [bubble]{сортування бульбашкою}, \hyperref [selection]{сортування вибором} та \hyperref [insertion]{сортування вставками}. І хоча в чистому вигляді ці алгоритми за рамками навчальних курсів практично не використовуються, вони важливі в теоретичному плані і лежать в основі багатьох ефективніших алгоритмів.

Під одномірним списком розумітимемо список, який містить інших списків чи кортежів. Як правило, завдання сортування ставиться для списків, елементи якого мають один і той самий тип.

Для певності припустимо, що ми маємо справу з цілим списком, який потрібно відсортувати за незменшенням.

\subsection{Сортування бульбашкою (bubble sort)}\label{bubble}

Суть алгоритму сортування бульбашкою \footnote{Також зустрічається інша назва "--- сортування обмінами.} полягає в послідовному попарному порівнянні двох сусідніх елементів списку. Якщо елементи розташовані в неправильному порядку (у нашому випадку "--- попередній більше наступного), вони змінюються місцями. Наприкінці одного такого проходу максимальний елемент спливе в кінці списку; метод отримав свою назву за аналогією зі спливаючими бульбашками вуглекислого газу в склянці з газованою водою. При наступному проході таким чином останній елемент можна не задіяти, і довжина невідсортованої частини списку зменшується на одиницю.

Тепер опишемо алгоритм трохи формальніше.
\begin{defn}{Алгоритм бульбашкового сортування (bubble sort)}
\begin{enumerate}
\item Робимо попарне порівняння сусідніх елементів. Якщо елемент \mi{[j]} більший за наступний елемент \mi{[j+1]}, міняємо їх місцями. Якщо на цьому етапі обмінів не відбувається, алгоритм завершується.
\item Повторюємо крок 1  для ще не відсортованих елементів.
\end{enumerate}
\end{defn}

Розглянемо спрощену програмну реалізацію алгоритму.

Зовнішній цикл реалізує другий крок алгоритму "---
повторення проходів. У найгіршому випадку таких проходів буде на одиницю менше довжини списку. Внутрішній цикл продає один прохід зліва направо. Зауважимо, що доцільність виконання наступного проходу за списком не перевіряється. Пропонуємо додати цей елемент самостійно.

\begin{py}{Сортування бульбашкою (bubble sort)}
# Введення цілого списку L
for i in range (len(L)-1) :
    for j in range (len(L)-1-i) :
        if L[j] > L[j+1] :
            tmp = L[j]
            L[j] = L[j+1]
            L[j+1] = tmp
# Висновок відсортованого списку
\end{py}

Обговоримо питання ефективності цього алгоритму з теоретичної точки зору. В теорії алгоритмів вводиться поняття складності алгоритму; якщо говорити спрощено, це оцінка зверху кількості операцій, необхідних для завершення алгоритму. При цьому використовується $ O $-символіка, прийнята у математичному аналізі.

Нагадаємо, що функція $ f (x)$ є $ O $ (велике) від $ g (x)$ при $ x \to \infty $, якщо знайдуться таке $ x _0$ і така константа $ C >0$, що для всіх $ x > x _0$ буде справедливою нерівність $\left | f (x) \right | \leqslant C \left | g (x) \right | $. Фраза \textit{<<алгоритм має складність $ O (f (n))$>>} означає, що \textit{зі збільшенням довжини $ n $ вхідної інформації кількість операцій (час роботи) алгоритму зростатиме не швидше, ніж деяка константа, помножена на $ f (n)$}.

\begin{defn}{Алгоритм вважається ефективним, якщо} \textbf{кількість операцій, необхідних для отримання відповіді у гіршому випадку, не дуже швидко зростає зі зростанням розміру вхідних даних.}
\end{defn}

У нашому випадку «розміром вхідних даних», очевидно, є довжина списку. Спробуємо обчислити кількість операцій, які будуть потрібні для завершення алгоритму в гіршому випадку.

Отже, якщо прийняти довжину списку, що дорівнює $ n $, то число проходів дорівнюватиме $(n -1)$ (ще раз підкреслимо, що тут принципово розглядається «гірший випадок»). Кожен прохід підсписку з $ k $ елементів вимагає $(k -1)$-го порівняння і (в~гіршому випадку!) $(k -1)$-го обміну, тобто в цілому $2(k -1)$ операції. Підсумкова кількість операцій складе
$$
\underbrace{2(n -1)+2(n -2)+\ldots +2}_{\text{всього}(n -1)\text{доданок}}=2\cdot \frac{n (n -1)}{2} = \left (n ^2- n \right) \text{операцій},
$$
т. е. складність даного алгоритму "--- $ O \left (n ^ 2 \right) $.

Недоліком цього алгоритму є велика кількість обмінів.

\subsection{Сортування вибором (selection sort)}\label{selection}

Ідея сортування вибором полягає у наступному. Знайдемо мінімальний елемент і переставимо його на початок списку (потрібно лише один обмін, на відміну від бульбашкового сортування). Далі знайдемо мінімальний елемент серед елементів, що залишилися, і переставимо його на друге місце. Продовжуючи процес $(n -1)$ разів, отримаємо відсортований список.

\begin{defn}{Алгоритм сортування вибором (selection sort)}
\begin{enumerate}
\item Знаходимо індекс мінімального елемента ще не відсортованої частини списку.
\item Здійснюємо обмін цього значення зі значенням у першій ще не відсортованій позиції (обмін не потрібен, якщо мінімальний елемент вже знаходиться на даній позиції).
\item 	Сортуємо <<хвіст>> списку, виключивши з розгляду вже відсортовані елементи.
\end{enumerate}
\end{defn}

\begin{py}{Сортування вибором (selection sort)}
# Введення цілого списку L
k = 0 # перша ще не відсортована позиція
while k < len(L)-1 :
    # Пошук позиції мінімального елемента
    min_pos = k
    for i in range (k+1, len(L)) :
        if L[i] < L[min_pos] :
            min_pos = i
    # Обмін
    tmp = L[k]
    L[k] = L[min_pos]
    L[min_pos] = tmp
    # Перехід до наступного кроку
    k = k + 1
# Висновок відсортованого списку
\end{py}

Наведена вище реалізація є спрощеною. Оцінимо кількість операцій даного алгоритму. Пошук мінімуму серед $ k $ елементів вимагає $(k -1)$-го порівняння; потім провадиться один обмін. Таким чином, кожен етап алгоритму вимагає $ k $ операцій, де $ k $ "--- довжина невідсортованої частини списку на даний момент. Усього маємо
$$
n + (n -1) + (n -2) + \ldots +1 = \frac{n (n +1)}{2} \text{операцій}.
$$
Незважаючи на те, що в середньому алгоритм сортування вибором вимагатиме менше операцій, ніж бульбашкова сортування, його складність також є $ O \left (n ^ 2 \right) $. По суті, це означає, що на списках досить великої довжини ми не відчуємо різниці між часом роботи обох алгоритмів.

\subsection{Сортування вставками (insertion sort)}\label{insertion}

Ідея сортування вставками полягає в тому, щоб вставити поточний елемент у правильне місце відсортованої частини списку. На початковому етапі відсортованим вважатимуться перший елемент; далі довжина відсортованої частини зростає на одиницю.

\begin{defn}{Алгоритм сортування вставками (insertion sort)}
\begin{enumerate}
\item Вважаємо $k=1$.
\item 	Нехай $ k $ елементів списку вже відсортовані. Беремо наступний елемент (\mi{L [k]}) і вставляємо його в потрібне місце відсортованої частини списку.
\item 	Збільшуємо $ k $ на одиницю і повторюємо крок 2. Алгоритм завершується, коли довжина невідсортованої 	частини дорівнює нулю.
\end{enumerate}
\end{defn}

\begin{py}{Сортування вставками (insertion sort)}
# Введення цілого списку L
for i in range (1, len (L)):
    key = L [i] # елемент, який необхідно вставити у потрібну позицію
    j = i -1
    while j >= 0 and L [j] > key : # шукаємо потрібну позицію справа наліво
        L [j +1] = L [j]
        j = j - 1
    L [j +1] = key # вставляємо
# Висновок відсортованого списку
\end{py}

Складність даного алгоритму також $ O \left (n ^ 2 \right) $ (перевірте це самостійно).

На \myhref{http :// airtucha. github. io / SortVis /}{даному ресурсі} представлена непогана візуалізація відомих алгоритмів сортування з відкритим кодом.

\subsection{Сортування складніших структур}

Як зазначалося, вбудованого методу \mi{sort()} буває недостатньо вирішення завдань. Нехай, наприклад, у списку \mi{L} зберігається деяка база даних:
\begin{py}{}
L = []
L.append (('Джонсон', (1, 1, 1980), 'бухгалтер', 800))
L.append (('Джексон', (15, 7, 1975), 'дантист', 2000))
L.append (('Андерсон', (8, 11, 1990), 'програміст', 1200))
L.append (('Стивенсон', (20, 5, 1992), 'касир', 400))
L.append (("Пітерсон", (11, 12, 1975), "адміністратор", 1000))
\end{py}
\noindent В даному випадку метод \mi{sort()} виконає сортування за першим елементом кортежу, тобто. на прізвища в алфавітному порядку. Але що, якщо нам необхідно відсортувати його за незменшенням заробітної плати (останній елемент кортежу)? Вирішимо це завдання за допомогою сортування вибором.

\begin{py}{}
k = 0
while k < len(L)-1 :
    min_pos = k
    for i in range (k+1, len(L)) :
        if L[i][3] < L[min_pos][3] : # порівнюються тільки зарплати
            min_pos = i
    tmp = L[k]
    L[k] = L[min_pos]
    L[min_pos] = tmp
    k = k + 1
# Виведення відсортованого списку: кожен елемент в окремому рядку
for x in L :
    print(x)
\end{py}

\section{Двовимірні списки (матриці)}

Вище ми вже бачили на прикладах, що потенціал списків у мові \Python\ аж ніяк не обмежується зручним поданням одновимірних (лінійних) масивів даних. У цьому параграфі розглянемо ще один найважливіший з практичної точки зору тип даних "--- \textbf{двовимірні таблиці}, або \textbf{матриці}. Як самостійний тип даних матриці в класичному \Python \відсутні, якщо не підключати додаткові бібліотеки\footnote{Див., наприклад , бібліотеку \myhref{http://www.numpy.org/}{NumPy}.}.

По суті, матриця є список рядків, кожен з яких також є список (наприклад, цілих чисел). При цьому взагалі довжини рядків не зобов'язані збігатися. У більш класичному випадку прийнято говорити про \textbf{прямокутну матрицю} розміром $ m \times n $ ($ m $ рядків і $ n $ стовпців), довжини всіх рядків якої однакові. Далі будемо припускати це за замовчуванням.

Звернення до елемента матриці $ M $, що стоїть у $ i $-му рядку в $ j $-му стовпці, має вигляд \mi{M[i][j]}. Конструкція \mi{M[i]} повертає список всіх елементів $ i $-го рядка матриці. Для отримання списку елементів стовпця потрібні складніші дії.

\subsection{Способи введення матриць}

\subsubsection{Перерахування елементів}

\begin{pyr}
\begin{py}{Приклад: створення матриці перерахуванням елементів}
M = [[1, 2, 3], [4, 5, 6]] # створюємо матрицю розміром 2х3
print (' M :', M)
\end{py}
\tcblower
\begin{py}{}
>>>
M: [[1, 2, 3], [4, 5, 6]]
>>>
\end{py}
\end{pyr}
\noindent Звернемо увагу на те, як виводиться матриця за допомогою стандартного \mi{print()}. Жодних відмінностей від виведення одномірного списку немає "--- всі елементи перераховуються через кому в квадратних дужках. Більш зручні способи виведення матриць будуть розглянуті \hyperref [moutput]{далі}.

\subsubsection{Використання конструктора}

Спробуємо створити ту саму матрицю за допомогою конструктора.

\begin{pyr}
\begin{py}{Приклад: створення матриці за допомогою конструктора списків}
M = [[3*i+j+1 for j in range(3)] for i in range (2)]
print('M:', M)
\end{py}
\tcblower
\begin{py}{}
>>>
M: [[1, 2, 3], [4, 5, 6]]
>>>
\end{py}
\end{pyr}
\noindent Цей код має сенс розібрати докладніше. В даному випадку ми використовували вкладений конструктор для створення одного рядка (з трьох елементів) та зовнішній конструктор для створення списку рядків (з двох елементів). Тут і далі індекс \mi{[i]} відповідає за номер рядка, індекс \mi{[j]} "--- за номер стовпця.

\subsubsection{Введення з клавіатури}

Матриці невеликого розміру у навчальних прикладах можна вводити з клавіатури. При цьому введення можна організувати або поелементно (натискаємо \keys{Enter} після кожного елемента), або рядково (вводиться рядок, елементи якого поділяються пробілом, а потім \keys{Enter}).

У першому випадку нам знадобляться два \mi{for}-цикли: зовнішній за номерами рядків і внутрішній за номерами стовпців. У внутрішньому циклі відбувається введення елементів \mi{x}, які збираються в рядок \mi{L}, який потім додається до матриці \mi{M}.

\begin{py}{Приклад: введення матриці з клавіатури поелементно}
m = int (input ('Кількість рядків:'))
n = int (input ('Кількість стовпців:'))
print ('Введіть елементи матриці:')
M = [] # порожній список для майбутньої матриці
for i in range (m) :
L = [] # перед введенням елементів рядка її потрібно обнулити
    for j in range (n) :
        x = int (input ()) # вводимо елемент
        L.append (x) # додаємо його до рядка
    M.append (L) # додаємо рядок до матриці
\end{py}

Цей код можна скоротити, якщо обидва цикли в конструктори списків. Подумайте, як це зробити.

При строковому введенні можна обійтися одним циклом. Черговий рядок будується на основі розбиття (\mi{split()}) введеної послідовності елементів. Зауважимо також, що цей спосіб не вимагає задання числа стовпців.

\begin{py}{Приклад: введення матриці з клавіатури по рядках}
m = int (input ('Кількість рядків:'))
print ('Введіть матрицю по рядках:')
M = []
for i in range (m) :
    s = input()
    L = [int(x) for x in s.split('')]
    M.append(L)
\end{py}

\subsubsection{Заповнення випадковими числами}

Тестування багатьох програм зручно виконувати на випадково згенерованих матрицях. Створимо випадкову цілочисленну матрицю розміром $ m \times n $.

\begin{py}{Приклад: заповнення матриці випадковими числами}
from random import *
m = int(input('Кількість рядків:'))
n = int (input ('Кількість стовпців:'))
M = []
for i in range (m) :
    L = []
    for j in range (n) :
        x = randint(-10, 10)
        L.append(x)
    M.append(L)
\end{py}
\noindent Використання конструкторів дозволяє мінімізувати код:
\begin{py}{Приклад: заповнення матриці випадковими числами}
from random import *
m = int(input('Кількість рядків:'))
n = int (input ('Кількість стовпців:'))
M = [[randint(-10, 10) for j in range(n)] for i in range(m)]
\end{py}

\subsubsection{Читання з файлу}

Часто необхідно обробити двовимірний масив даних, що зберігається в електронній таблиці. Приклад нижче показує, як прочитати дані з \textsf{CSV}-файлу в цілу матрицю. Передбачається, що комірки електронної таблиці містять лише числа, без жодних описів та заголовків.

\begin{py}{Приклад: читання матриці з файлу}
f = open('m.csv', 'r')
M = []
for line in f :
    L = [int(x) for x in line.split(';')]
    M.append(L)
f.close()
\end{py}

\subsection{Виведення матриць у <<красивому>> вигляді}\label{moutput}

Щоб вивести матрицю на екран у зручнішій для читання формі, іноді досить вивести її рядково:
\begin{pyr}
\begin{py}{Приклад: виведення матриці строково}
M = [[1, 2, 3], [4, 5, 6]]
for row in M :
    print(row)
\end{py}
\tcblower
\begin{py}{}
>>>
[1, 2, 3]
[4, 5, 6]
>>>
\end{py}
\end{pyr}
\noindent Результат не ідеальний, але все ж таки краще, ніж стандартний \mi{print()}. Спробуємо прибрати квадратні дужки та коми:
\begin{pyr}
\begin{py}{Приклад: виведення матриці у <<красивому>> вигляді}
M = [[1, 2, 3], [4, 5, 6]]
for i in range(len(M)) : 
    for j in range (len(M[i])) : 
        print(M[i][j], end=' ') #  виводимо елементи рядка через пробіл
    print() #  перехід на новий рядок
\end{py}
\tcblower
\begin{py}{}
>>>
1 2 3
4 5 6
>>>
\end{py}
\end{pyr}
\noindent Тут \mi{len(M)} "--- число рядків (довжина матриці), \mi{len(M[i])} "--- число стовпців (довжина рядка матриці). Такий варіант виведення виглядає краще попереднього, однак і він виглядатиме кострубато у випадку, наприклад, матриці з випадковими елементами з діапазону $ [-100; 100] $:
\begin{py}{}
>>>
-95 -68 98
55 3 40
>>>
\end{py}
Нормальне сприйняття результату заважає різна довжина елементів, а також знаки «мінус» у деяких з них. Тому має сенс вдосконалити наше виведення на екран, вирівнюючи числа по лівому чи правому краю. Тут нам несподівано знадобиться один з методів \hyperref[rjust]{\mi{ljust()}, \mi{rjust()}} для рядків. Нагадаємо, ці методи дозволяють доповнити рядок пробілами до потрібної довжини з необхідним вирівнюванням. Традиційно числові дані вирівнюються праворуч:
\begin{pyr}
\begin{py}{Приклад: виведення матриці у <<красивому>> вигляді з вирівнюванням}
from random import *
M = [[randint(-100, 100) for j in range(6)] for i in range(4)]
for i in range (4) :
    for j in range (6) :
        print(str(M[i][j]).rjust(4), end=' ')
    print()
\end{py}
\tcblower
\begin{py}{}
>>>
15 -71 94 -16 -23 -8
36 29 -86 -43 52 8
47 66 -90 -23 26 -88
49 16 -86 -76 -12 -11
>>>
\end{py}
\end{pyr}

Ось тепер "--- те що треба! Параметр \mi{4} для \mi{rjust()} обраний тому, що максимальна довжина виведеного числа могла скласти 4 символи (число \mi{-100}).

Цикл виведення матриці на екран можна скоротити до одного рядка коду, якщо застосувати конструктор і метод \mi{join()}. Спробуйте зробити це самостійно.

Використання різних методів обробки рядків дозволяє налаштовувати за власним бажанням вирівнювання елементів матриці, роздільники рядків та стовпців, а також інші параметри виведення.

\section{Завдання}

У цьому параграфі акцент зроблено на завданнях навчального характеру, що показують основні прийоми роботи з одновимірними та двовимірними масивами даних з використанням списків та матриць.

\begin{pr}{}
Створити список малих літер англійської абетки.
\end{pr}

Для вирішення скористаємося функціями \mi{ord()} та \mi{chr()}.

\begin{py}{Розв'язання}
a = ord('a')
z = ord('z')
alphabet = [chr(x) for x in range (a, z+1)]
print(alphabet)
\end{py}

\begin{pr}{}
Створити список перших 20 трицифрових натуральних чисел із сумою цифр, кратною 7.
\end{pr}

Ідея Розв'язання у наступному. Спочатку створюємо порожній список \mi{Nums}, куди будемо поміщати знайдені елементи. Основний цикл працює до тих пір, поки довжина цього списку не стане рівною 20. У самому циклі відбувається наступне: знаходимо суму цифр чергового числа \mi{i} і перевіряємо її подільність на 7. У разі позитивної відповіді записуємо число \mi{i} у список \mi{Nums}. Зрозуміло, перед початком циклу потрібно задати початкове значення \mi{i}. В даному випадку беремо \mi{100}, оскільки розглядаються лише тризначні числа. Для продовження циклу \mi{i} збільшуємо на одиницю.

\begin{py}{Розв'язання}
Nums = []
i = 100
while len(Nums) < 20 :
    s = i % 10 + (i // 10) % 10 + i // 100
    if s % 7 == 0 :
        Nums.append(i)
    i = i + 1
print(Nums)
\end{py}

\begin{pr}{}
У заданому списку списку знайти мінімальний елемент і його позицію.
\end{pr}

Застосовується типовий алгоритм пошуку найбільших і найменших значень послідовності. Зверніть увагу, що нам не обов'язково заводити змінну для зберігання мінімального значення "--- достатньо зберігати лише номер позиції \mi{min_pos}. Тоді до самого значення можна отримати доступ за індексом: \mi{L [min_pos]}.

\begin{py}{Розв'язання}
# Введення списку L
min_pos = 0
for i in range (len(L)) :
    if L[i] < L[min_pos] :
        min _ pos = i
print ('Мінімальний елемент дорівнює', L [min _ pos], \
'і знаходиться в позиції', min _ pos)
\end{py}

Зауважимо, що в результаті ми отримаємо позицію лише першого входження мінімального значення до списку.

При розв'язанні цієї задачі були потрібні позиції елементів, тому був застосований розширений варіант \mi{for}-циклу для перегляду списку:
\begin{defn}{Перегляд списку із запам'ятовуванням позицій елементів}
\begin{py}{}
for i in range (len(L)) : # i - індекс списку
    # Обробка елементів; L [i] - елемент списку
\end{py}
\end{defn}

\begin{pr}{}
У заданому списку списку знайти кількість парних елементів, що стоять на непарних місцях.
\end{pr}

Перед нами типове завдання одновимірного пошуку елементів із заданими властивостями. Опишемо алгоритм розв'язання. Спочатку кількість \mi{k} необхідних елементів покладемо рівним нулю. Далі у циклі переглядаємо список поелементно. Якщо поточний елемент задовольняє заданій властивості (умові), то змінну \mi{k} збільшуємо на одиницю. Після закінчення циклу виводимо значення змінної \mi{k} на екран. Подібним чином обчислюються і суми елементів із заданою властивістю.

\begin{py}{Розв'язання}
# Введення списку L
k = 0
for i in range (len(L)) :
    if L[i] % 2 == 0 and i % 2 == 0 :
        k = k + 1
print (k)
\end{py}
\noindent У нашому випадку умова відбору включає як перевірку самих значень \mi{L[i]}, так і їх індексів \mi{[i]}. Звертаємо увагу, що <<непарні місця>> інтерпретуються як парні індекси з огляду на те, що нумерація індексів елементів списку ведеться з нуля, а <<місць>> - з одиниці.

\begin{pr}{}
Дано цілочисельний список. Чи вірно, що у ньому немає негативних елементів?
\end{pr}

Для вирішення завдання нам знову доведеться аналізувати значення елементів списку, але при цьому вже не враховуватимуться індекси. Отже, можемо скористатися коротким варіантом \mi{for}-циклу для перегляду списку:
\begin{defn}{Перегляд списку: короткий варіант циклу}
\begin{py}{}
for x in L :
    # Обробка елементів; x - елемент списку
\end{py}
\end{defn}

\begin{py}{Розв'язання}
# Введення списку L
for x in L :
    if x < 0 :
        print('No')
        break
else :
    print('Yes')
\end{py}

Зазначимо, що негативну відповідь на наше запитання можна отримати, зустрівши перше негативне значення. Для ствердної відповіді доведеться пройти весь список до кінця. Тут зручно використовувати конструкцію \mi{for / else}.

\begin{pr}{}
У заданому списку списку видалити всі елементи, що стоять після непарних.
\end{pr}

Завдання, що вимагають додавання або видалення елементів списку, містять пастку. Щоб не траплятися до неї, запам'ятайте правило:
\begin{defn}{Якщо довжина списку може змінюватися в процесі,}
то для перегляду елементів не можна застосовувати \mi{for}-цикли.
\end{defn}

\noindent Пов'язано це з тим, що в \mi{for}-циклі діапазон індексів елементів списку задається a ~ priori у заголовку циклу і не може бути змінений. Але якщо довжина списку зміниться в процесі перегляду, то ми або не дійдемо до кінця списку (у разі додавання елементів), або індекс вийде за межі списку (у разі видалення). Виходом тут може бути створення нового списку для зберігання результату, проте подібні дії зазвичай не вітаються через неекономне використання пам'яті.

Проблему вирішує \mi{while}-цикл, умова закінчення якого перевіряється багаторазово. Це дозволить коректно змінювати довжину списку всередині циклу як більшу, так і меншу сторону.

\begin{defn}{\texttt{while}-цикл для перегляду списку}
\begin{py}{}
i = 0
while i < len(L) :
    # Обробка елементів; L[i] - елемент списку
    i = i + 1
\end{py}
\end{defn}

При вирішенні цього завдання нам доведеться переглядати список з кінця, оскільки при перегляді від початку в ситуації, коли два непарні елементи знаходяться підряд (\mi{L[i]} і \mi{L[i +1]}), буде видалено елемент \mi{L[i+1]}, і, таким чином, наступний за ними елемент,\mi{L[i+2]}, може бути пропущений (а за умовою завдання \mi{L[i+2]} також має бути видалено).

\begin{py}{Розв'язання}
# Введення списку L
i = len (L)
while i > 0 : #перегляд тільки до другого елемента
    if L [i -1] % 2 == 1 : # якщо попередній елемент непарний
        del L [i] # видаляємо поточний
    i = i - 1 # зсуваємо індекс поточної позиції на початок
# Виведення списку L
\end{py}

\begin{pr}{}
У заданому списку цілих чисел вставити 0 після кожного негативного елементу.
\end{pr}

Тут також основний цикл не може бути \mi{for}-циклом, тому що довжина списку, можливо, буде змінюватися. Вставку нульових елементів виконаємо методом \mi{insert()}.

\begin{py}{Розв'язання}
# Введення списку L
i = 0
while i < len(L):
    if L [i] < 0: # якщо поточний елемент негативний
        L.insert (i +1, 0) # вставляємо 0 у позицію після поточного
        i = i + 2 # "перестрибуємо" через 0
    else:
        i = i + 1 # а в цьому випадку стрибати не потрібно
# Виведення списку L
\end{py}

\begin{pr}{}
Заповнити матрицю розміром $ n \times n $ за наступним правилом:
$$
\begin{pmatrix}
1 & 0 & 0 & \ldots & 0 \\
2 & 2 & 0 & \ldots & 0 \\
3 & 3 & 3 & \ldots & 0 \\
\vdots & \vdots & \vdots & \ddots & \vdots \\
n&n&n&\ldots&n
\end{pmatrix}.
$$
\end{pr}

Наше завдання "--- виявити закономірність у заповненні матриці числами і скласти формулу. Вся верхня половина (над головною діагоналлю) матриці заповнена нулями. Індекси елементів цієї частини матриці підпорядковуються правилу: номер стовпця $ j $ строго більше за номер рядка $ i $ (перевірте це).Отже, для інших елементів матриці вірно протилежна умова.Далі, елементи кожного рядка, що стоять до головної діагоналі включно, рівні номеру цього рядка (нагадаємо, номер і індекс "--- не зовсім одне і те ж). Тепер можна написати код.

\begin{py}{Розв'язання}
n = int(input(' n = '))
M = []
for i in range(n):
L = [] # список-заготівля для майбутнього рядка матриці
    for j in range(n) :
        if j > i : # якщо елемент у "верхній" частині матриці
            L.append(0) # додаємо 0
        else :
            L.append(i+1) # інакше - додаємо номер рядка
    M.append(L) # додаємо рядок до матриці
# Виведення матриці на екран
\end{py}

\begin{pr}{}
У заданій цілочисельній матриці знайти суму негативних елементів.
\end{pr}

Пошук елементів у матриці за заданою властивістю потребує обходу всіх елементів. Так як матриця "--- двомірна структура, то знадобиться два цикли (зовнішній по рядках і вкладений за елементами рядків).

\begin{py}{Розв'язання}
# Введення матриці M
S = 0
for i in range (len(M)) :
    for j in range (len(M[i])) :
        if M[i][j] < 0 :
            S = S + M[i][j]
print ('Сума негативних елементів:', S)
\end{py}
\noindent Замість \mi{len(M)} і \mi{len(M [i))} можна використовувати змінні \mi{m} і \mi{n} (кількість рядків і стовпців), якщо вони були задані явно або введено з клавіатури.

\begin{defn}{Цикл для перегляду елементів матриці}
\begin{py}{}
for i in range (len(M)) :
    for j in range (len(M[i])) :
        # обробка елементів M [i][j]
\end{py}
\end{defn}

У цьому можна було застосувати і короткий варіант циклу, оскільки індекси елементів не знадобилися.

\begin{defn}{Цикл для перегляду елементів матриці: короткий варіант}
\begin{py}{}
for row in M :
    for x in row :
        # обробка елементів x рядка row
\end{py}
\end{defn}

\begin{pr}{}
Дана цілочисленна квадратна матриця. Чи правильно, що сума елементів головної діагоналі поділяється на три?
\end{pr}

Головна діагональ матриці характеризується властивістю $ i = j $, де $ i $"--- номер рядка, $ j $"--- номер стовпця. Отже, для перегляду всіх елементів головної діагоналі можна обійтися одним циклом "--- $ i $ (або $ j $, що не має значення).

Параметр \mi{len (M)} у разі квадратної матриці є як число рядків, так і число стовпців.

\begin{py}{Розв'язання}
# Введення матриці M
S = 0
for i in range (len(M)) :
    S = S + M[i][i]
if S % 3 == 0 :
    print('Yes')
else :
    print('No')
\end{py}

\begin{pr}{}
Дана цілочисленна матриця. Чи правда, що хоча б один із рядків починається з нуля?
\end{pr}

Для вирішення цього завдання нам знадобляться лише перші елементи кожного рядка, тому цикл буде один "--- за номерами рядків".

\begin{py}{Розв'язання}
# Введення матриці M
for i in range (len(M)) :
    if M [i] [0] == 0: # M [i] [0] - перший елемент i-го рядка
        print('Yes')
        break
    else :
        print('No')
\end{py}
\noindent Як тільки знайдено рядок, що починається з нульового елемента, можна завершувати цикл і відповідати <<Ні>>. Відповісти <<Так>> ми зможемо лише після завершення перегляду всіх рядків.

\begin{pr}{}
У заданій цілочисельній матриці видалити рядки з парною сумою елементів.
\end{pr}

Як і в разі одномірних списків, якщо довжина матриці може змінитися в процесі обробки, для перегляду елементів не можна використовувати \mi{for}-цикл. Зауважимо, проте, що у цій задачі може змінюватися лише число рядків (тобто \mi{len (M)}, довжина матриці), але з число стовпців. Тому для обчислення суми елементів рядка \mi{for}-цикл цілком підходить.

\begin{py}{Розв'язання}
# Введення матриці M
i = 0 # номер поточного рядка
while i < len (M) :
    S = 0 # обнуляємо суму елементів на початку кожного рядка
    for j in range(len(M[i])) :
        S = S + M[i][j]
    if S % 2 == 0 : # якщо сума елементів парна
        del M[i] # видаляємо рядок M[i]
    else :
        i = i + 1 # інакше - переходимо до наступного рядка
# Виведення матриці M
\end{py}
\noindent Зверніть увагу: збільшувати лічильник \mi{i} не потрібно, якщо рядок було видалено: у цьому випадку поточне значення \mi{i} вже буде вказувати на наступний рядок матриці.

\begin{pr}{}
У заданій цілочисельній матриці видалити стовпці, що містять хоча б один негативний елемент.
\end{pr}

З видаленням стовпців у матриці справа складніша. Якщо рядок \mi{[i0]} можна видалити однією командою: \mi{del M[i0]}, то для стовпців у нас немає такого зручного позначення. По суті, видалення стовпця \mi{[j0]} означає, що ми повинні видалити елементи з індексом \mi{[j0]} для кожного рядка.

Для вирішення задачі зробимо таке:
\begin{enumerate}
\item При першому перегляді матриці \mi{M} створимо список \mi{Cols_to_del} (columns to delete, стовпці для видалення) номерів стовпців, що містять негативні елементи.
\item Створимо нову матрицю \mi{New_M}. Переглядаючи вихідну матрицю \mi{M} повторно, запишемо в нову матрицю елементи, індекси стовпців яких не належать списку \mi{Cols_to_del}.
\end{enumerate}

\begin{py}{Розв'язання}
# Введення матриці M
Cols_to_del = []
for j in range (len (M [0])) : # цикл за номерами стовпців
    for i in range (len (M)) : # цикл для перегляду елементів стовпця
        if M[i][j] < 0 :
            Cols_to_del.append(j)
            break
New_M = []
for i in range (len(M)) :
    L = []
    for j in range(len(M[i])) :
        if j not в Cols_to_del :
            L.append(M[i][j])
    New_M.append(L)
# Виведення матриці New_M
\end{py}
\noindent Варто звернути увагу на те, як організовано перегляд матриці \mi{M} у першому пункті. Зовнішній та внутрішній цикли переставлені місцями для того, щоб переглядати елементи по стовпцях, а не по рядках. Як тільки негативний елемент знайдено, номер поточного стовпця заноситься до «чорного списку» і цикл перегляду елементів стовпця закінчується.

При такому порядку перегляду в якості числа стовпців у матриці було використано довжину першого рядка: \mi{len(M[0])}\footnote{Подумайте, чому тут не можна було написати стандартне \mi{len(M[i])}.}. Це нічого не суперечить, адже ми спочатку припускаємо, що довжини всіх рядків однакові.

\begin{pr}{}
Знайти добуток двох заданих матриць.
\end{pr}

Нагадаємо, добуток $C=A\cdot B$ двох матриць $A$ і $B$ визначено, якщо число стовпців матриці $A$ дорівнює числу рядків матриці $B$. Якщо матриця $A$ має розмірність $m\times n$, а матриця $B$ "--- $n\times p$, то розмірність результуючої матриці $C$ є $m\times p$. Елементи матриці $C$ обчислюються за формулою
$$
c_{ij}=\sum\limits_{k=1}^{n}a_{ik}b_{kj}, \qquad i=\overline{1, m}, \quad j=\overline{1, p}.
$$
де $a_{ik}$ і $b_{kj}$ "--- елементи матриць $A$ і $B$ відповідно. Операція множення матриць некомутативна.

\begin{py}{Розв'язання}
# Введення матриць A та B
if len(A[0]) != len(B) :
    print('Помилка. Твір не визначено')
else :
    C = [] # починаємо обчислювати матрицю C
    for i in range (len(A)) :
        L = [] # починаємо обчислювати рядок матриці C
        for j in range (len(B[0])) :
            s = 0 # починаємо обчислювати елемент C[i][j]
            for k in range (len(B)) :
                s = s + A [i] [k] * B [k] [j]
            L.append(s)
        C.append(L)
# Виведення матриці C
\end{py}


% ідея завдання - Решето Ератосфена за допомогою списків генераторів
